\subsection{Gamification}
Studien innefattar inte någon praktisk instans av gamification utan reflekterar snarare kring dess tillämpning i en C-DAS-OB-applikation och vad som är viktigt kopplat till gamification ur lokförarens perspektiv.
Detta i enlighet med studiens syfte, omfattning och forskningsfrågor.
När gamification ska tillämpas i C-DAS-applikationer är det viktigt att det inte görs på ett sätt så att lokförarna inte känner sig utvärderade utan istället med hjälp av gamification kan utvärdera sig själva.
Tillämpandet av gamification och intrinsisk motivation utanför spelsvärden är en relativt ny trend och sett till det förvånansvärt passionerade engagemang dataspel frambringar kan potentialen mycket väl vara stor vid rätt utförande.
Gamification är dock komplext och kräver djuplodande kunskaper i såväl beteendevetenskap som datavetenskap och inte minst domänkunskap för slutanvändare och det system varuti det är tänkt att implementeras.
